\section{Pseudoscience and Scientism}

\emph{In which we review several books by Wolfgang Smith on the metaphysics of science.}

\begin{quotex}
It is impossible that these things should be understood by men in general, but only by the small number of those who are destined to prepare in one way or another the germs of the future cycle.
\flright{\textsc{Rene Guenon}}
\end{quotex}


\textbf{Boris Mouravieff} identified three stages in the intellectual development of the West: philosophy, theology, and science. The coming stage, the age of Art, will be the creative synthesis, the contemplation of Beauty in mysticism and theurgy as described by \textbf{Vladimir Solovyov} in the Law of Development\footnote{\url{https://gornahoor.net/?p=3571}}.

The aim of philosophy is the good life or the best regime, the aim of theology is knowledge of God, and that of science, knowledge of the physical world. These correspond to the Taoist triad of Man, Heaven, and Earth. In an essay titled \textit{Modern Science and the Guenonian Critique}, \textbf{Wolfgang Smith} challenges Guenon's casual dismissal of profane science. Imagine two intersecting circles of a Venn diagram. Beyond the intersection of one circle there is ``pseudoscience'', beyond the other is ``scientism''. In the intersection is what Smith calls ``hard science''. Smith admits that scientism, i.e., the reduction of everything to quantity, is not knowledge, yet he insists that there is a partial truth in ``hard science''. Hence, it may be helpful to consider Smith's work as a ``Guenonian critique of modern science.''

We agrees with this, since those preparing for the next cycle should be adept in political science, metaphysics, and, yes, science. All the more so since the best minds today go into science, although they seldom have any grasp of metaphysics. Moreover, we would go further than Smith, since the creative act requires a full integration. \textbf{Valentin Tomberg} explains:

\begin{quotex}
The synthesis of science and religion is not a theory, but rather the inner act of consciousness of adding the spiritual vertical to the scientific horizontal.
\end{quotex}


Due to the large amount of material, it would be impracticable to go through each book in detail. It may be more useful to those unfamiliar with Smith's works to outline his fundamental worldview, since it is scattered across several essays. Then we can look at some specific examples from physics, cosmology, biology, and psychology that illustrate how profane science can and should be evaluated from a higher, metaphysical viewpoint. This may make his books, which are far from casual reading, more approachable.

\paragraph{Principles of Traditional Cosmology}


In the \emph{Wisdom of Ancient Cosmology} (WAC), Smith identifies four principles to which all ancient cosmology conforms.

\begin{itemize}


\item Traditional cosmology primarily deals with the qualitative aspects of reality.

\item The metaphysical notion of verticality affirms a hierarchic order in which the corporeal domain constitutes the lowest tier.

\item Man constitutes a microcosm or ``universe in miniature'', which recapitulates the cosmic order. Thus man is not a stranger in a ``hostile or indifferent universe but constitutes the very heart and center of the cosmos''.

\item The higher strata of the cosmos can be known through the realization of the corresponding states of man himself. Thus the key to knowledge is first to ``know thyself''.
\end{itemize}


\paragraph{Bifurcation}


Smith traces the foundations of science, as we know it today, to Rene Descartes and Galileo Galilei. The introduced some errors into the consciousness of modern minds:

\begin{itemize}


\item Bifurcation

\item Quantification

\item De-essentialism
\end{itemize}


Bifurcation is Smith's word for Cartesian dualism: all that exists is matter and mind. Matter is mechanistic and the objects of sense experience exist only in the mind. The physical world, therefore, is knowable only as quantity so that mathematics is the only tool to know it. All qualities, then, exist only in mind. That is why Guenon asserts that

\begin{quotex}
In our world, by reason of the special conditions of existence to which it is subject, the lowest point takes on the aspect of pure quantity, deprived of every qualitative distinction.
\end{quotex}


The net result is the de-essentialization of thought. An essence or form is the answer to the ``what is'' question. Scientism denies the existence of such essences. Philosophically, that is the doctrine of ``nominalism'', which asserts that ideas are merely human creations, as opposed to ``realism'', which claims that the ideas are real and independent of man.

Although we cannot accept their respective philosophical systems \emph{in toto}, Smith points to \textbf{Alfred North Whitehead} and to \textbf{Edmund Husserl} as the two philosophers of the 20th century who did the most to provide compelling alternatives to Cartesian bifurcation.

\paragraph{Vertical Causation}


The ideas of vertical and horizontal causation are fundamental to Smith's critique. Nevertheless, they are poorly understood, even by those who consider themselves religious today. Science is concerned solely with horizontal causes, that is, events that follow each other in time, in other words, with efficient causality. From a methodological perspective, it omits any consideration of vertical causation, that is, formal and final causes. Obviously, science has had great success, at least in terms of technology, with that method. However, scientism goes on to assert, without any warrant, that vertical causation does not exist at all.

Another way to look at this question, which Smith hints at without being so explicit, is the Hermetic teaching of the triad of Providence, Will, and Destiny\footnote{\url{https://gornahoor.net/?p=8116}}. Providence, then, includes formal and final causes. Specifically, formal causes determine the possibilities of manifestation of those forms that are realizable at any moment. The final cause is the ultimate aim or purpose of the manifest universe.

Ontologically the final cause is prior to the efficient cause, even if the efficient cause appears to come first in time. Hence, scientifically a final cause is expressed as one of the anthropic principles, which will be discussed later. Smith also refers to \textbf{Edgar Dacqué}'s idea that man was created prior to the animals, so that the process of evolution is part of that creation in time. He uses the same argument used by \textbf{Julius Evola} in \textit{The Esoteric Origin of the Species}\footnote{See Section \ref{sec:InvolutionEvolution} in this book.}.

Scientists and atheists often react to the deformed ideas of God and vertical causation held by naïve believers. This view seems itself to be bifurcationist. It envisions a ``spiritual'' world, inhabited by a god, that runs along in time parallel to the physical world. From time to time, this god looks down and interferes with the world process. This notion is commonly derided as the ``god of the gaps''.

Of course, God is outside of time altogether. Smith refers to the Scholastic idea of \emph{nunc stans} or the ``eternal now'' to express that view. This does not mean that time is a sequence of ``nows'', just as a line is not ``made'' of points. Hence, horizontal causation, or Destiny, can best be understood as a projection of vertical causation onto time.

\paragraph{Metaphysics as Seeing}


Smith asserts that metaphysics ``springs from man's innate thirst for truth, which is none other than the thirst for God''. Today, it is often restricted to its degenerate form as an academic discipline pursued by ``professionals''. Rather, metaphysics is an activity of the mind and heart to which we are ``called''. Smith claims, democratically, that in principle all are called; from our perspective---and Guenon's for that matter---that is far from obvious.

Gnosis then is not a matter of debate or arguing, is more like ``seeing'', or a direct perception. The purpose of logic is to ``deconstruct false beliefs'' and perforce to \emph{purify the mind}. Only the pure in heart shall see God.

Drawing on Husserl's phenomenology, Smith claims that the phenomenon is that which ``shows itself in itself''. This ``seeing'' of the self-revelation of the phenomenon is non-dual; the bifurcation into the seeing subject and the seen object comes afterward. If we mistake the bifurcation for the real, then ``seeing'' has come to an end. Although Smith comes at this from a completely different angle, his use of phenomenology corresponds to what we have called Hermetic meditation\footnote{\url{https://gornahoor.net/?p=898}}.

By this process, we understand the object to be a ``whole'' in itself, not merely the sum of its parts. The horizontal causation described by science necessarily misses that point in its attempt to reduce everything to more fundamental parts. For example, the organism is its genes, an object is ``made'' of atoms, atoms, of subatomic particle, and so on. The scientist thereby can only ``see'' his own theories and constructs and misses the actual phenomenon. This is far from the ``knowledge of the whole'' which is the proper aim of philosophy.

For support, Smith turns to \textbf{Johann Wolfgang von Goethe} who conceived of a science not of quantities but of qualities. Goethe claimed that the mind must be ``sun-like'' in order to behold the Sun. Smith goes further, based on some ideas from \textbf{Meister Eckhart}, that the mind must be divine in order to behold God. That is the ultimate knowledge of the whole.

Part 2 of this review will deal with specific scientific topics discussed by Smith. These will be taken from the following books:

\begin{itemize}
\item \textit{The Wisdom of Ancient Cosmology: Contemporary Science in Light of Tradition} (WAC) \item \textit{Science and Myth} (S\&M) \item \textit{Cosmos \& Transcendence: Breaking Through the Barrier of Scientistic Belief} (C\&T) \item \textit{The Quantum Enigma: Finding the Hidden Key} (QE) 
\end{itemize}

On a personal note, I found my copy of QE at a used book store, taken apparently from Huston Smith's personal library. There is this hand-written dedication:

\begin{quotex}
To Professor Huston Smith in friendship, gratitude and high esteem, Wolfgang
\end{quotex}


The book was in pristine condition and seemed, to me, to have never been read.


\flrightit{Posted on 2016-05-16 by Cologero}

\begin{center}* * *\end{center}

\begin{footnotesize}
\begin{sffamily}
\texttt{Tom Walker on 2016-05-17 at 09:12 said:}

One of the best books debunking the new atheism and scientistic pretentions is that of David Berlinski '' The Devil's Delusion'' , \url{http://www.davidberlinski.org/devils-delusion/about.php}

Berlinski , a secular Jew , appears to be neutral as far as religion is concerned but in his book he dismantles atheistic and scientistic arrogance and shows that science , is far from having an answer for everything even in the physical world .

\hfill

\texttt{Cologero on 2016-05-18 at 20:25 said:}

Thanks for the link, Tom Walker, but our intent is far from merely ``debunking the new atheism and scientistic pretentions''. Rather, as I thought was made clear, is to take what is valid and place it into a larger metaphysical context, while avoiding both pseudoscience (which is bad science masquerading as philosophy) and scientism (which is bad philosophy masquerading as science).

\hfill

\texttt{Sigurd on 2017-09-07 at 12:43 said:}

Wolfgang Smith's works are great. I was introduced to them from this website. It's nice to read a more contemporary writer on the subject of the same issues, particularly one who is already a specialist, Keith Critchlow would be another good example.

On the subject of cosmology, Cologero do you think theurgy and ceremonial ritual employed in the context of planetary correspondences, such as is suggested in Evola's hermetic tradition, is a valid practice in contemporary times? I understand there is some psychic danger in this because of the ``objectifying'' of the various forces that correspond to the padmas, what are your thoughts?

\hfill

\texttt{littleM on 2018-05-16 at 09:18 said:}

I always stand in awe, thinking how everything here is written before it happened. Well, I should get used to it, but that does not mean being less in awe.

\hfill

\texttt{N. W. Flitcraft on 2018-05-16 at 20:59 said:}

If you are not already, you should all be aware of Dr. Smith's new foundation, the \textit{Philos-Sophia Initiative}\footnote{\url{https://philos-sophia.org/}}. And follow the \textit{Philos-Sophia Initiative ``page'' on facebook}\footnote{\url{https://www.facebook.com/PhilosSophiaInitiative/}} and/or join in the conversation at the \textit{Wolfgang Smith Facebook Circle}\footnote{\url{https://www.facebook.com/groups/Wolfgangsters/}}.

\end{sffamily}
\end{footnotesize}

